% =====================================================================
%  Chapter 00 — How to read this book
% =====================================================================
% The master's `record' environment was malformed (\colorbox's argument
% brace was split across \newenvironment's begin- and end-code, so TeX
% absorbed the end-code into the begin-code).  Diagnosed here at S48 and
% REPAIRED IN THE MASTER PREAMBLE via \newsavebox/lrbox; the local
% \renewenvironment that stood here is therefore removed as redundant.

\chapter{How to Read This Book}
\label{ch:how-to-read}

This book is a record of a research programme that did not achieve its goal.
It is written for a mathematician who has never seen the programme, and its
purpose is to make the whole of it --- the reasoning, the constructions, the
numerical apparatus, the corrections and the retractions --- legible and
checkable by somebody who was not present for any of it.

Three things must be said before anything else.

\begin{record}
\textbf{First.} This book \emph{states} results. It does not \emph{prove}
them. The proofs, where they exist, live in
\texttt{proofs/claim\_a\_3d\_proof\_attempt.tex}, a \(10{,}462\)-line
\LaTeX{} document that compiles to \(142\) pages. Every mathematical
statement in this book is a report on an object in that document, cited by
its label and, where useful, its line number.

\smallskip
\textbf{Second.} The programme has \emph{not} proved the Clay Millennium
Prize problem. It claimed a proof once, in session~\Sn{29}, and withdrew that
claim in session~\Sn{31} on eight independent grounds. Global regularity for
general \(u_0\in H^1(\T^3)\) is \OPEN{} in the proof document, and is
recorded as such throughout.

\smallskip
\textbf{Third.} The programme's credibility, such as it is, rests entirely on
the failures being reported as crisply as the successes. A wall is a wall. A
withdrawn claim is withdrawn. Nothing in the chapters that follow softens a
negative result, and where a negative result was once softened, the softening
is itself reported.
\end{record}

\section{The status labels}
\label{sec:status-labels}

Every substantive claim in this book carries exactly one status label. The
labels are not decoration; they are the compressed verdict, and they come
from the classification scheme the proof document itself adopted in
\texttt{sec:walls-s47} (line~9717).

\begin{center}
\small
\begin{tabular}{@{}l p{0.72\textwidth}@{}}
\toprule
Label & Meaning \\
\midrule
\PROVED
  & A proof exists in the proof document, is unconditional within its stated
    hypothesis class, and has survived every audit since. \\[2pt]
\OPEN
  & Nothing is proved against it. It is unfinished and remains available. An
    \OPEN{} item is not a weak \PROVED{}; it is an absence of a proof. \\[2pt]
\IMPOSSIBLE
  & A proof in the document rules the route out, and no choice of parameters
    inside the route can help. This label is used only where such a proof
    exists, and it is never weakened in the prose. \\[2pt]
\WITHDRAWN
  & An assertion was made in the document and has been retracted. The
    retraction is permanent unless the underlying defect is repaired and
    named as repaired. \\[2pt]
\SUPERSEDED
  & The route was replaced by a different formulation of the same target,
    which carries its own status. Superseded is not a synonym for refuted. \\[2pt]
\PROVEDMOD{X}
  & Proved except for enumerated verifications of class~\(X\); the missing
    items involve no new estimate. \\[2pt]
\CONDITIONAL{(H)}
  & Proved on the assumption of a named hypothesis \((H)\) which is itself
    not proved here. The hypothesis must be carried in the same sentence as
    the conclusion. \\
\bottomrule
\end{tabular}
\end{center}

Two of these distinctions did real work and are worth dwelling on.

\begin{obsv}[\IMPOSSIBLE{} is narrower than it sounds]
\label{obs:impossible-narrow}
Three of the programme's ten walls carry \IMPOSSIBLE{}, and each is
\emph{qualified}. Wall~4 (nested level iteration) is proved impossible for
that route, while the target it was meant to serve is merely \SUPERSEDED{}.
Wall~8 is proved impossible \emph{as written}, with a hypothetical localized
substitute left \OPEN{}. Wall~9 is proved impossible \emph{as an import}, and
the underlying question it concerns is of course \OPEN{} --- it is the Prize.
Where a qualification exists, this book carries it.
\end{obsv}

\begin{obsv}[The most important label in the book is \OPEN]
\label{obs:open-is-honest}
The programme's central structural obstruction --- the Calder\'on--Zygmund
logarithm \(\Lfac\) of Chapter~\ref{ch:inheritance} --- is \OPEN{}, not
\IMPOSSIBLE{}. No proof in the document says the logarithm cannot be removed.
What is established is that it is paid once, at a single identifiable step,
and inherited everywhere downstream. Recording it as a proved barrier would
have overstated it; recording it as a minor bookkeeping nuisance would have
understated it. It is neither.
\end{obsv}

\section{Session tags}
\label{sec:session-tags}

The programme ran as forty-seven working sessions between April and September
2026. A tag of the form \Sn{31} names session thirty-one. Tags are used
constantly and mean something precise: they locate a claim in time, and they
identify the session log in \texttt{SESSION-LOG/} that carries the primary
account of what happened. Session \Sn{31} is the audit; \Sn{29} is the Prize
claim it withdrew; \Sn{40} is the comparison with the forced-blowup
construction of Chapter~\ref{ch:openai}; \Sn{46} is the strategic audit and
\Sn{47} the walls document.

Sessions are not uniform in size or kind. Some are a single analytical
derivation, some are a numerical production run, one (\Sn{31}) is an audit
that produced no new mathematics and changed everything. Numbering has gaps
where a session produced nothing worth logging; \Sn{39} carries two logs
because two independent tracks ran under it.

\begin{center}
\small
\begin{tabular}{@{}l l p{0.52\textwidth}@{}}
\toprule
Range & Period & Character \\
\midrule
\Sn{1}--\Sn{17} & Apr 2026 & The numerical programme: property engine, 2D
  solvers, 3D solvers, diagnostics. Milestones M0--M10. \\
\Sn{18}--\Sn{27} & May 2026 & The Claim-A proof attempt
  (Chapter~\ref{ch:claim-a}). \\
\Sn{29}--\Sn{30} & May 2026 & The Prize claim and the \(\sigma\)-framework
  (Chapter~\ref{ch:prize-claim}). \\
\Sn{31} & Jul 2026 & The audit; withdrawal (Chapter~\ref{ch:audit}). \\
\Sn{32}--\Sn{39} & Jul 2026 & The alignment pincer and profile rigidity. \\
\Sn{40} & Sep 2026 & The forced-blowup comparison
  (Chapter~\ref{ch:openai}). \\
\Sn{41}--\Sn{45} & Sep 2026 & The annulus residual, the band conjecture, the
  scale diagnostics. \\
\Sn{46}--\Sn{47} & Sep 2026 & Strategic audit; first successful compile; the
  walls document. \\
\bottomrule
\end{tabular}
\end{center}

\section{The source hierarchy}
\label{sec:source-hierarchy}

The programme generated four kinds of record, and they do not always agree.
When they disagree, the following order decides, and the disagreement is
footnoted rather than silently resolved.

\begin{figure}[ht]
\centering
\begin{tikzpicture}[
  node distance=6mm,
  box/.style={draw, rounded corners=2pt, align=center, inner sep=4pt,
              font=\small, minimum width=52mm},
  lbl/.style={font=\scriptsize\itshape, align=left, text width=52mm}
]
\node[box, fill=panelgrey] (doc)
  {\textbf{1. The proof document}\\[1pt]
   \texttt{proofs/claim\_a\_3d\_proof\_attempt.tex}};
\node[lbl, right=4mm of doc] (doclbl)
  {governs all mathematical statements; compiles; labels are primary, line
   numbers secondary};

\node[box, below=of doc] (db)
  {\textbf{2. The results database}\\[1pt]
   \texttt{results/results.db} (31 tables)};
\node[lbl, right=4mm of db] (dblbl)
  {governs all quantitative experimental claims: \texttt{claim\_id},
   \texttt{key\_metric}, \texttt{verdict}};

\node[box, below=of db] (logs)
  {\textbf{3. The session logs}\\[1pt]
   \texttt{SESSION-LOG/} (40 files)};
\node[lbl, right=4mm of logs] (logslbl)
  {govern narrative: what was tried, why, what was believed at the time};

\node[box, below=of logs] (prog)
  {\textbf{4. The summaries}\\[1pt]
   \texttt{PROGRESS.md}, \texttt{HANDOVER-S46-OPUS.md}};
\node[lbl, right=4mm of prog] (proglbl)
  {navigational only; known to contain errors of omission; never decisive};

\draw[-{Latex[length=2mm]}, thick] (doc) -- (db);
\draw[-{Latex[length=2mm]}, thick] (db) -- (logs);
\draw[-{Latex[length=2mm]}, thick] (logs) -- (prog);

\node[left=3mm of doc, font=\scriptsize, rotate=90, anchor=south]
  {decreasing authority};
\end{tikzpicture}
\caption[The source hierarchy]{The source hierarchy. Where two sources
disagree, the higher one governs and the disagreement is footnoted. This is
not a hypothetical policy: applying it in \Sn{47} caught three overstatements
in the summaries, one of them material.}
\label{fig:source-hierarchy}
\end{figure}

The rule earned its place. In \Sn{47}, checking the summaries against the
document found three claims the document does not support. The most
consequential: \texttt{PROGRESS.md} records the \Sn{31} audit as having
\emph{seven} grounds. The document's \texttt{sec:audit-s31} (line~5218) lists
\emph{eight}, D1--D8. The omitted one, D8, is the first appearance of the
Calder\'on--Zygmund logarithm --- the obstruction that came to dominate every
later session. Losing it from the summary was not neutral: it made the
programme's central difficulty look like a later regression rather than
something the audit had already named.\footnote{The two other \Sn{47}
corrections: \texttt{HANDOVER-S46-OPUS.md} described
\texttt{conj:lc-s44} as ``a common \emph{weaker} target'' than
\texttt{conj:K-refined}, whereas by \texttt{rem:s44-lc-reading} (line~9453)
neither implies the other; and the phrase ``\(c_K\) and band-absorption
proved independent'' drops the qualifier that \texttt{rem:s44-scope}
(line~8822) insists on --- the counterexamples are over admissible test
configurations, not over Navier--Stokes solutions.}

The reader should assume more such errors remain in the summaries. This book
was written against the document, the database and the logs, in that order.

\section{Citing the proof document}
\label{sec:citing}

References into the proof document are given as text, never as a \LaTeX{}
cross-reference: for example \texttt{prop:transfer-audit}, line~6941. Labels
are primary and stable; line numbers are secondary, quoted as of \Sn{47}, and
will rot as the document grows. Where this book gives a line number it has
been checked against the current file.

One trap is pre-existing and worth stating once. Parts of \S20--\S22 of the
proof document carry \emph{hand-written} numbers in their prose and
subsection headings --- ``Lemma 20.1'', ``Theorem 20.4'' --- which are off by
one from the numbers \LaTeX{} assigns; those statements live in the section
\LaTeX{} numbers 21. The drift was exposed by the document's first successful
compilation, in \Sn{47}, and has not been repaired. This book uses labels, so
it is unaffected, but a reader who follows a hand-written number will be sent
to the wrong section.

\section{A minimal notation crib}
\label{sec:crib}

The full glossary is an appendix. These few symbols recur from
Chapter~\ref{ch:origin} onwards and are worth fixing now.

\begin{center}
\small
\begin{tabular}{@{}l l p{0.60\textwidth}@{}}
\toprule
Symbol & Reads as & Meaning \\
\midrule
\(T(x,t)\) & --- & Temperature: a \emph{scalar} field, the master variable of
  the founding thesis (Chapter~\ref{ch:origin}). \\
\(A(T)\) & --- & Thermal diffusivity \(k(T)/(\rho c_v)\), required strictly
  positive by the second law. \\
\(\theta\) & theta & The auxiliary scalar of Route~2:
  \(\partial_t\theta+u\cdot\nabla\theta=\nu\Delta\theta+\nu\abs{\nabla u}^2\). \\
\(M(t)\) & --- & \(\norm{\omega(\cdot,t)}_{L^\infty}\), the vorticity maximum. \\
\(\rstar\) & r-star & The self-similar scale \(M^{-1/2}\). \\
\(\sstar\) & sigma-star & The instantaneous coherence deficit, \texttt{def:sigma-star}
  (line~5393): a scale-invariant measure of vortex-direction disorder at
  scale \(\rstar\). \\
\(\Lfac\) & script-L & \(\log(e+\norm{u}_{H^3}/M)\), the Calder\'on--Zygmund
  logarithm. The single most important symbol in this book. \\
\(\ehat\) & e-hat & The vorticity direction field \(\omega/\abs{\omega}\). \\
\(\Aann\) & --- & The annulus residual of \Sn{36}--\Sn{41}. \\
\bottomrule
\end{tabular}
\end{center}

\section{The shape of the book}
\label{sec:shape}

Part~I gives the founding thesis, the three routes it generated, the
architecture it planned, and the verification discipline it eventually
learned. Part~II is the numerical programme, layer by layer, with the actual
numbers. Part~III is \emph{the early work}: the Claim-A proof attempt, the
Prize claim that grew out of it, the audit that withdrew it, and --- the
chapter that carries the most weight for the rest of the book --- what the
later work inherited from all three. Part~IV is the later analytical
programme. Part~V places the programme against a contemporaneous result on
the forced equations. Part~VI is the walls. Part~VII is what the programme
learned about method, and Part~VIII is where it stands.

The early work is retained in full and read forward. It was not deleted, and
it is not quietly superseded. That decision was the owner's, and it is the
right one: the programme's single most consequential structural finding was
first written down in the audit that destroyed its central claim, and a
history that discarded the failure would have discarded the finding with it.
